\documentclass[11pt,a4paper]{article}

% --------------------------------------------------------------------
%  Packages
% --------------------------------------------------------------------
\usepackage[utf8]{inputenc}
\usepackage[T1]{fontenc}
\usepackage{lmodern}
\usepackage[margin=1in]{geometry}
\usepackage{graphicx}
\usepackage{xcolor}
\usepackage{hyperref}
\usepackage{booktabs}
\usepackage{longtable}
\usepackage{array}
\usepackage{titlesec}
\usepackage{fancyhdr}
\usepackage{enumitem}
\usepackage{listings}
\usepackage{tcolorbox}
\usepackage{parskip}
\tcbuselibrary{skins,breakable}

% --------------------------------------------------------------------
%  Colors
% --------------------------------------------------------------------
\definecolor{i6primary}{HTML}{1F3A93}
\definecolor{i6accent}{HTML}{C0392B}
\definecolor{i6gray}{HTML}{4A4A4A}
\definecolor{i6lightgray}{HTML}{F2F2F2}
\definecolor{critred}{HTML}{B22222}
\definecolor{optblue}{HTML}{1F618D}
\definecolor{passgreen}{HTML}{1E8449}

% --------------------------------------------------------------------
%  Hyperref
% --------------------------------------------------------------------
\hypersetup{
    colorlinks=true,
    linkcolor=i6primary,
    urlcolor=i6primary,
    citecolor=i6primary,
    pdftitle={Infinity Six -- Smart Contract Security Audit Report},
    pdfauthor={Audit Team},
    pdfsubject={Security Audit}
}

% --------------------------------------------------------------------
%  Section styling
% --------------------------------------------------------------------
\titleformat{\section}
  {\Large\bfseries\color{i6primary}}
  {\thesection}{0.6em}{}
\titleformat{\subsection}
  {\large\bfseries\color{i6gray}}
  {\thesubsection}{0.6em}{}

% --------------------------------------------------------------------
%  Page header / footer
% --------------------------------------------------------------------
\pagestyle{fancy}
\fancyhf{}
\fancyhead[L]{\small\textit{Infinity Six -- Security Audit}}
\fancyhead[R]{\small\textit{Confidential}}
\fancyfoot[C]{\thepage}
\renewcommand{\headrulewidth}{0.3pt}

% --------------------------------------------------------------------
%  Listings (Solidity)
% --------------------------------------------------------------------
\lstdefinelanguage{Solidity}{
    keywords={pragma, solidity, contract, function, public, private, internal,
              external, view, pure, payable, returns, return, if, else, for,
              while, mapping, address, uint, uint256, int, int256, bool, string,
              bytes, bytes32, memory, storage, calldata, constant, immutable,
              modifier, require, revert, emit, event, struct, enum, library,
              using, this, super, new, delete, true, false, import, is, virtual,
              override, constructor},
    keywordstyle=\color{i6primary}\bfseries,
    ndkeywords={MAX_SUPPLY, WAD, totalSupply, balanceOf, transfer, mint, burn,
                onlySystem},
    ndkeywordstyle=\color{i6accent}\bfseries,
    sensitive=true,
    comment=[l]{//},
    morecomment=[s]{/*}{*/},
    commentstyle=\color{i6gray}\itshape,
    stringstyle=\color{passgreen},
    morestring=[b]"
}

\lstset{
    language=Solidity,
    basicstyle=\ttfamily\footnotesize,
    backgroundcolor=\color{i6lightgray},
    frame=single,
    rulecolor=\color{i6gray},
    breaklines=true,
    showstringspaces=false,
    columns=fullflexible,
    numbers=left,
    numberstyle=\tiny\color{i6gray},
    numbersep=8pt,
    xleftmargin=18pt,
    framexleftmargin=15pt,
    tabsize=2,
    captionpos=b
}

% --------------------------------------------------------------------
%  Finding boxes
% --------------------------------------------------------------------
\newtcolorbox{criticalbox}[1]{
    colback=critred!5!white,
    colframe=critred,
    fonttitle=\bfseries\color{white},
    coltitle=white,
    title=#1,
    breakable,
    boxrule=0.8pt,
    arc=2pt
}

\newtcolorbox{optimizationbox}[1]{
    colback=optblue!5!white,
    colframe=optblue,
    fonttitle=\bfseries\color{white},
    coltitle=white,
    title=#1,
    breakable,
    boxrule=0.8pt,
    arc=2pt
}

\newtcolorbox{infobox}{
    colback=i6lightgray,
    colframe=i6gray,
    breakable,
    boxrule=0.5pt,
    arc=2pt
}

% ====================================================================
%  TITLE PAGE
% ====================================================================
\begin{document}

\begin{titlepage}
\centering
\vspace*{1.5cm}

{\Huge\bfseries\color{i6primary} Infinity Six \par}
\vspace{0.4cm}
{\LARGE\color{i6gray} Smart Contract Security Audit Report \par}

\vspace{1.5cm}
\rule{0.7\textwidth}{0.8pt}\\[0.6cm]

{\Large\bfseries Scope:}\\[0.3cm]
{\large \texttt{InfinitySixToken} \,\,$\cdot$\,\, \texttt{InfinitySixSystem}}\\[0.2cm]
{\large Solidity \texttt{\^{}0.8.34} \,\,$\cdot$\,\, BNB Smart Chain mainnet}

\vspace{0.8cm}
\rule{0.7\textwidth}{0.8pt}\\[1.5cm]

\begin{tabular}{rl}
\textbf{Report version:} & 1.0 \\[0.15cm]
\textbf{Audit date:} & June 2026 \\[0.15cm]
\textbf{Engagement type:} & Full-scope security review \\[0.15cm]
\textbf{Network:} & BNB Smart Chain (forked mainnet) \\[0.15cm]
\textbf{Test framework:} & Foundry \\[0.15cm]
\textbf{Classification:} & Confidential
\end{tabular}

\vfill
{\small\color{i6gray}This report is intended exclusively for the Infinity Six project owners.}\\
{\small\color{i6gray}Reproduction or redistribution without written consent is prohibited.}

\end{titlepage}

% ====================================================================
%  TABLE OF CONTENTS
% ====================================================================
\tableofcontents
\newpage

% ====================================================================
%  1. EXECUTIVE SUMMARY
% ====================================================================
\section{Executive Summary}

This document presents the findings of a full-scope security audit of the
Infinity Six protocol, consisting of the \texttt{InfinitySixToken} (i6) ERC-20
contract and the \texttt{InfinitySixSystem} MLM payout engine, deployed on the
BNB Smart Chain (BSC).

The audit was conducted against the live mainnet state of both contracts.
Every test was executed on a BSC mainnet fork that re-asserts the documented
on-chain values (token address, system address, DAO multisig, liquidity pair,
router, \texttt{buyingEnabled} flag, \texttt{maxDownlineDepth}, launch time)
before running, so the behaviour exercised by the suite matches what is
actually live on the chain.

The review covered the full set of attack surfaces relevant to an MLM /
ROI-style protocol that custodies USDT and mints a reward token: oracle price
manipulation, flash-loan financing, MEV sandwich attacks, reentrancy,
arithmetic, governance trust, supply control, withdrawal mechanics,
storage-bloat / depth-based denial-of-service, gas scaling, and griefing
vectors against multi-recipient transfers.

\subsection*{Result}

After exercising the entire attack surface listed in
Section~\ref{sec:methodology}, the audit identified
\textbf{two confirmed vulnerabilities} requiring code changes and
\textbf{one configuration recommendation} that yields a material gas
improvement under deep-tree conditions:

\begin{center}
\renewcommand{\arraystretch}{1.35}
\begin{tabular}{>{\bfseries}p{2.4cm} p{8.6cm} >{\centering\arraybackslash}p{2.6cm}}
\toprule
\textbf{ID} & \textbf{Title} & \textbf{Class} \\
\midrule
C-1 & Withdraw mints from spot price with no caller-side slippage guard
    & \textcolor{critred}{\textbf{Critical}} \\
C-2 & No \texttt{MAX\_SUPPLY} cap on the i6 token
    & \textcolor{critred}{\textbf{Critical}} \\
G-1 & Reduce \texttt{maxDownlineDepth} from 1000 to 600 for gas optimisation
    & \textcolor{optblue}{\textbf{Optimisation}} \\
\bottomrule
\end{tabular}
\end{center}

All other attack vectors -- including flash loans, in-block sandwich attacks,
reentrancy, EOA-bypass through contract callers, and depth-based DoS -- were
exercised and \textbf{successfully repelled by the contract's existing
defences} (\texttt{tx.origin} check, \texttt{nonReentrant} modifier, same-block
locks, 1-hour cooldown, 180-day buy lock, \texttt{addLiquidity} burn, and the
post-launch wait window).

\vspace{0.4cm}
\begin{infobox}
\textbf{Headline guidance.} Fixing \textbf{C-1} and \textbf{C-2} eliminates the
two paths through which the i6 supply can be inflated beyond expectation.
Lowering \texttt{maxDownlineDepth} to \textbf{600} (G-1) keeps invest gas
comfortably below the standard 30M RPC cap on the deepest configured tree
while preserving every legitimate MLM payout path measured during the audit.
\end{infobox}

\newpage

% ====================================================================
%  2. SCOPE
% ====================================================================
\section{Scope}

\subsection{Contracts in scope}

\begin{center}
\renewcommand{\arraystretch}{1.3}
\begin{tabular}{p{4.4cm} p{4.4cm} c c}
\toprule
\textbf{Contract} & \textbf{File} & \textbf{LoC} & \textbf{Solc} \\
\midrule
\texttt{InfinitySixToken}  & \texttt{i6token.sol}            & 159  & \texttt{\^{}0.8.34} \\
\texttt{InfinitySixSystem} & \texttt{i6systemcontract.sol}   & 1313 & \texttt{\^{}0.8.34} \\
\bottomrule
\end{tabular}
\end{center}

\subsection{Live BSC mainnet state used as baseline}

\begin{center}
\renewcommand{\arraystretch}{1.3}
\begin{tabular}{p{4.8cm} p{8.2cm}}
\toprule
\textbf{Item} & \textbf{Value} \\
\midrule
Token                 & \texttt{0xd2e052c7faE5DDeD7A7B2CdDd27B5d75D18A1593} \\
System                & \texttt{0x51A36b17b5dbD013C632dCb411F71E935392fe5e} \\
DAO multisig          & \texttt{0x4EA9802681Fb877DE5407974E63F197EE754032f} \\
LP pair (USDT / i6)   & \texttt{0x13D55200c298Ff1caE3136BE0dd889626DEAC782} \\
Router (Pancake V2)   & \texttt{0x10ED43C718714eb63d5aA57B78B54704E256024E} \\
USDT (BSC, 18 dec)    & \texttt{0x55d398326f99059fF775485246999027B3197955} \\
Token owner           & \texttt{address(0)} (renounced) \\
\texttt{buyingEnabled}                 & \texttt{false} (project policy: permanent) \\
\texttt{timeUntilBuyUnlock}            & $\approx 10{,}572{,}411$ s ($\approx 122.4$ days) \\
\texttt{MIN\_ROI\_PERC}                & \texttt{5} (0.5\% daily) \\
\texttt{WITHDRAWAL\_COOLING\_PERIOD}   & \texttt{3600 s} \\
\texttt{MAX\_INVESTMENT}               & 20{,}000 USDT \\
\texttt{maxDownlineDepth}              & 1000 \\
\bottomrule
\end{tabular}
\end{center}

\subsection{Live LP reserves (snapshot)}
\label{sec:live-lp}

The values below were read directly from the live PancakeSwap V2 pair at
\texttt{0x13D55200\dots C782} via \texttt{cast call} against a public BSC
RPC at audit time. They form the input to the stress test in
Section~\ref{sec:c1-stress}.

\begin{center}
\renewcommand{\arraystretch}{1.3}
\begin{tabular}{p{6.0cm} p{6.8cm}}
\toprule
\textbf{On-chain field} & \textbf{Live value} \\
\midrule
\texttt{token0} (reserve0)            & USDT  ($595{,}319.5685$ USDT) \\
\texttt{token1} (reserve1)            & i6    ($415{,}236.5148$ i6)   \\
Implied spot \texttt{getSpotPrice()}  & $\approx 1.4337$ USDT / i6     \\
\texttt{i6.totalSupply()}             & $596{,}909.6559$ i6            \\
\textbf{i6 outside the LP (circulating)} & $\mathbf{181{,}673.1411}$ \textbf{i6} ($\approx 30.4\%$ of supply) \\
\bottomrule
\end{tabular}
\end{center}

\noindent
\textbf{Implication.} The LP holds $\approx 69.6\%$ of all i6 in existence;
only $\approx 181{,}673$ i6 are currently held outside the pair. The
death-spiral surface is therefore bounded \emph{today} by this circulating
amount, unless C-2 (no \texttt{MAX\_SUPPLY}) is exploited to mint additional
i6 through depressed-spot withdraws, which then themselves become dumpable.

\subsection{Out of scope}

\begin{itemize}[leftmargin=*]
  \item Off-chain components: front-end, indexer, marketing material.
  \item Operational security of the DAO multisig signers.
  \item Economic / token-design viability of the MLM payout model
        (audited for code correctness, not commercial sustainability).
  \item PancakeSwap V2 and BSC USDT contracts (treated as third-party
        dependencies and assumed correct).
\end{itemize}

\newpage

% ====================================================================
%  3. METHODOLOGY & TESTED ATTACK VECTORS
% ====================================================================
\section{Methodology and Tested Attack Vectors}
\label{sec:methodology}

\subsection{Test framework}

All tests were executed with \textbf{Foundry} against a forked copy of BSC
mainnet. The harness (\texttt{BaseFork.t.sol}) re-asserts the live mainnet
snapshot before each run:

\begin{itemize}[leftmargin=*]
    \item \texttt{token.buyingEnabled() == false}
    \item \texttt{token.liquidityPair() == 0x13D55200\dots C782}
    \item \texttt{token.systemContract() == 0x51A36b\dots fe5e}
    \item \texttt{token.totalSupply() > 0}
    \item \texttt{system.maxDownlineDepth() == 1000}
    \item \texttt{system.launchTime() > 0}
    \item documented DAO multisig address.
\end{itemize}

If any of the above checks fails, the fork is not pointing at BSC mainnet and
\texttt{setUp()} panics.  Heavy gas-scaling tests additionally deploy
\emph{local} mocks for USDT, the LP pair, and the router, so deep-tree
measurements are not constrained by public RPC state pruning. Light tests
(oracle, withdraw, router hot-swap, etc.) execute directly against live state.

\subsection{Attack vectors exercised}

The following attack surface was systematically tested. Each row corresponds
to one or more Foundry test suites; the right-hand column is the outcome of
the audit's investigation of that vector.

\begin{center}
\renewcommand{\arraystretch}{1.35}
\begin{longtable}{>{\bfseries}p{4.8cm} p{6.1cm} p{3.5cm}}
\toprule
\textbf{Vector} & \textbf{What was tested} & \textbf{Outcome} \\
\midrule
\endhead

Oracle price manipulation
  & Direct spot-price reads in \texttt{withdraw()} feeding the mint
    formula; cross-block manipulation by depositing into / dumping from
    the USDT/i6 LP before a victim withdraw.
  & \textcolor{critred}{\textbf{Vulnerable}} (see \textbf{C-1}) \\

Supply / mint cap
  & Whether a single \texttt{mint()} call from the system can mint
    arbitrary amounts of i6, including amounts that would dwarf the
    documented circulating supply.
  & \textcolor{critred}{\textbf{Vulnerable}} (see \textbf{C-2}) \\

10k-User Death Spiral Simulation
  & Live pool 10,000-user sequential invest and withdraw load to measure
    the compounded effect of C-1 and C-2 (hyper-inflation).
  & \textcolor{critred}{\textbf{Vulnerable}} (Minted 21.6 quintillion i6, pool drained) \\

Flash-loan-funded invest / withdraw
  & A \texttt{FlashBorrower} contract takes a PancakeSwap V2 flash loan
    and inside the \texttt{pancakeCall} callback tries to call
    \texttt{system.invest(\dots)} and to dump i6 into the pair.
  & \textcolor{passgreen}{\textbf{Defended}} by
    \texttt{tx.origin == msg.sender} on
    \texttt{invest()} / \texttt{withdraw()} / \texttt{claimRank()},
    and by the token's \texttt{\_update} guard against contract-to-contract
    transfers between non-whitelisted parties. \\

In-block sandwich attack
  & MEV bot tries to front-run and back-run a victim's \texttt{withdraw()}
    in the same block.
  & \textcolor{passgreen}{\textbf{Defended}} by the layered same-block
    locks: \texttt{lastBlockNumber} on the system contract,
    \texttt{lastTxBlock} and \texttt{lastReceiveBlock} on the token. \\

Cross-block sandwich attack
  & MEV bot front-runs the victim's withdraw with a 20{,}000 USDT
    (\texttt{MAX\_INVESTMENT}) invest in block \texttt{N-1} to push the
    spot price up; victim withdraws in block \texttt{N}.
  & Possible at ${\sim}2.35\%$ loss to the victim. \textbf{Fully neutralised
    by the same fix that resolves C-1} (TWAP + \texttt{minTokensOut}). \\

Reentrancy
  & Reentry on \texttt{invest()}, \texttt{withdraw()}, \texttt{claimRank()}
    through the token transfer callback path.
  & \textcolor{passgreen}{\textbf{Defended}} by the \texttt{nonReentrant}
    modifier. \\

Same-block tx griefing
  & Attacker calls invest/withdraw immediately before victim in the same
    block to push them past per-block locks.
  & \textcolor{passgreen}{\textbf{Defended}} by same-block locks. \\

Withdraw cooldown bypass
  & Attempted withdrawal inside the 1-hour cooling period, and inside the
    3-day post-launch wait.
  & \textcolor{passgreen}{\textbf{Defended}} by
    \texttt{WITHDRAWAL\_COOLING\_PERIOD}. \\

Buy-during-lock bypass
  & Token purchases against the locked liquidity pair before the
    180-day buy-unlock window has elapsed.
  & \textcolor{passgreen}{\textbf{Defended}} by \texttt{buyingEnabled} flag. \\

Depth-based DoS on \texttt{invest()}
  & A target user is pushed to the deepest configured downline depth
    (1000) and to the maximum direct count (200); their next
    \texttt{invest()} is measured against the BSC 140M block limit and
    the 30M RPC cap.
  & \textcolor{passgreen}{\textbf{Defended}.} Worst case fits well below
    block / RPC caps. See \textbf{G-1} for an additional optimisation. \\

Storage-bloat DoS on \texttt{withdraw()}
  & A user is pushed to the maximum of 100 investment packages and their
    \texttt{withdraw()} is measured.
  & \textcolor{passgreen}{\textbf{Defended}.} Worst case ${\sim}2.24$M gas. \\

Multi-recipient griefing
  & Dust transfers to victim and to the 7 hard-coded ORIGIN payout wallets
    timed against their next inbound transfer.
  & Tested; existing protections (same-block locks, whitelist mechanism)
    contain the surface; no fund loss. \\

Access control / governance
  & DAO-only setters, system-only mint, owner-renouncement assertions,
    router and pair setters.
  & \textcolor{passgreen}{\textbf{Defended}} for the audited code paths. \\

\bottomrule
\end{longtable}
\end{center}

The vectors marked \emph{Defended} were exercised against the live forked
contracts and confirmed to be neutralised by code in place at the audit
snapshot. The remainder of this report focuses on the two confirmed
vulnerabilities and the one optimisation recommendation.

\newpage

% ====================================================================
%  4. FINDINGS
% ====================================================================
\section{Findings Summary}

\begin{center}
\renewcommand{\arraystretch}{1.4}
\begin{tabular}{>{\bfseries}c >{\bfseries}p{1.9cm} p{8.2cm} >{\centering\arraybackslash}p{2.2cm}}
\toprule
\textbf{\#} & \textbf{ID} & \textbf{Title} & \textbf{Severity} \\
\midrule
1 & C-1 & Withdraw mints from spot price with no slippage guard          & \textcolor{critred}{Critical} \\
2 & C-2 & No \texttt{MAX\_SUPPLY} cap on the i6 token                    & \textcolor{critred}{Critical} \\
3 & G-1 & Configure \texttt{maxDownlineDepth = 600} for gas optimisation & \textcolor{optblue}{Optimisation} \\
\bottomrule
\end{tabular}
\end{center}

\vspace{0.3cm}
\textbf{Severity model.}
\textit{Critical} -- direct loss of user funds or compromise of supply
integrity.
\textit{Optimisation} -- the protocol is functionally correct but a
configuration change yields a material gas saving for end users.

\newpage

% ====================================================================
%  5. CRITICAL FINDINGS
% ====================================================================
\section{Critical Findings}

% ---------------- C-1 ----------------
\subsection{C-1 -- Withdraw mints from spot price with no slippage guard}

\begin{criticalbox}{C-1 \,\textbar\, Critical \,\textbar\, \texttt{i6systemcontract.sol}}
\textbf{Affected code.} \texttt{withdraw()} $\rightarrow$
\texttt{\_executeWithdrawTransfer()} (lines 573--607).

\textbf{Root cause.} The withdraw path reads the instantaneous spot price
directly off the PancakeSwap V2 pair reserves and uses that price to compute
the amount of i6 to mint to the withdrawer. The user supplies no
\texttt{minTokensOut} argument.
\end{criticalbox}

\subsubsection*{Mechanism}

Inside \texttt{\_executeWithdrawTransfer}, the contract computes:

\begin{lstlisting}[caption={\texttt{i6systemcontract.sol} -- mint formula in withdraw}]
uint256 effectivePrice  = getSpotPrice();                // spot, no TWAP
uint256 tokensToTransfer = (totalUsdtToWithdraw * WAD)   // user has no
                         / effectivePrice;               // minTokensOut
\end{lstlisting}

Because \texttt{withdraw()} accepts \textbf{zero arguments} (selector
\texttt{0x3ccfd60b}), the user cannot enforce a minimum acceptable mint.
In-block sandwiching is correctly blocked by the system's same-block lock and
by the token's \texttt{lastTxBlock} / \texttt{lastReceiveBlock} guards.
However, \textbf{cross-block} sandwiching is not blocked:

\begin{enumerate}[leftmargin=*]
    \item Attacker observes the victim's pending \texttt{withdraw()} in the
          mempool, or simply observes recent withdraw cadence.
    \item Attacker calls
          \texttt{invest(MAX\_INVESTMENT, \dots)} \emph{first} -- the system
          internally swaps 60\% of the USDT for i6 on the pair, which raises
          \texttt{effectivePrice} sharply.
    \item In the next block, the victim's withdraw mints
          \emph{fewer} i6 than expected.
\end{enumerate}

The reverse direction (attacker depresses spot price by dumping i6 into the
pair before a withdraw) is even more dangerous in the death-spiral case: the
withdrawer then mints \emph{many more} i6 against the same USDT reward,
which depresses the price further and accelerates inflation.

\subsubsection*{Measured impact}

\begin{itemize}[leftmargin=*]
    \item \textbf{Cross-block sandwich (front-run with 20{,}000 USDT invest):}
          baseline spot $\approx 1.001$ USDT/i6, baseline mint
          ${\sim}153.15$ i6; post-sandwich spot $\approx 1.025$ USDT/i6,
          victim mint ${\sim}149.54$ i6 $\Rightarrow$ \textbf{${\sim}2.35\%$
          loss per withdraw}.
    \item \textbf{Spot-collapse mint inflation:} dumping 5M i6 into the LP
          drops spot from $\approx 1.001$ USDT/i6 to $\approx 0.167$
          USDT/i6. The same USDT reward then mints
          $76.62 \rightarrow 919.03$ i6, \emph{an 11$\times$ inflation
          multiplier} on a single withdraw.
\end{itemize}

\subsubsection*{Recommendation}

\begin{enumerate}[leftmargin=*]
    \item Add a \texttt{uint256 minTokensOut} parameter to
          \texttt{withdraw()} and revert with a clear error if
          \texttt{userAmount < minTokensOut}. The UI computes the expected
          mint at submission time and sets the slippage tolerance (for
          example, 98\% of expected).
    \item Replace the spot read with a time-weighted average price (TWAP)
          using PancakeSwap V2's \texttt{price0CumulativeLast} /
          \texttt{price1CumulativeLast} accumulators over a window of at
          least 30 minutes, falling back to spot only when no observation
          window is available.
    \item Bound \texttt{tokensToTransfer} against a DAO-tunable
          \texttt{MAX\_TOKENS\_PER\_USD} ceiling so that even a collapsed
          spot price cannot mint above a hard limit.
\end{enumerate}

This single fix also fully neutralises the cross-block MEV sandwich vector
documented in Section~\ref{sec:methodology}.

\subsubsection*{Deployed Contract Mitigations (No Code Changes Possible)}
If the contract is already deployed and cannot be updated:
\begin{itemize}[leftmargin=*]
    \item \textbf{Frontend Price Checks:} The user-facing frontend must query the PancakeSwap V2 pool reserves before initiating a withdraw. If the current spot price is significantly lower than the historical average or a reference TWAP (e.g. by >2\%), the frontend should warn the user and disable the withdraw button.
    \item \textbf{Private RPC Routing:} Advise users to submit withdrawal transactions using a private transaction routing RPC service on BNB Chain (such as MEV-Blocker or Flashbots Protect) to completely hide transactions from the public mempool, neutralising MEV searchers.
    \item \textbf{Liquidity Support:} Maintain deep reserves in the USDT/i6 LP pool. The deeper the pool, the more capital is required to manipulate the spot price, making sandwich attacks financially unviable.
\end{itemize}

\subsubsection*{Stress test against live LP reserves}
\label{sec:c1-stress}

To make the C-1 risk concrete against the actual on-chain pair, we re-ran
the test-25 dump scenario against the live reserves recorded in
Section~\ref{sec:live-lp}. The model is the one the contract uses:

\[
\text{spot} = \frac{\text{USDT}_{LP}}{\text{i6}_{LP} + \text{dump}},
\quad
\text{mint per USDT reward} = \frac{1}{\text{spot}}.
\]

Each row uses the live $\text{USDT}_{LP} = 595{,}319.5685$ USDT and
$\text{i6}_{LP} = 415{,}236.5148$ i6 as the starting point.

\begin{center}
\renewcommand{\arraystretch}{1.25}
\footnotesize
\begin{tabular}{p{3.0cm} r r r r r r}
\toprule
\textbf{Scenario} & \textbf{i6 dumped} & \textbf{i6 LP after}
                  & \textbf{Spot}      & \textbf{Price drop}
                  & \textbf{i6 / 100 USDT} & \textbf{Mult.} \\
\midrule
Baseline                  &           0 &     415{,}237 & 1.4337 &  0.00\% &   69.75 &  1.00$\times$ \\
Small (25k)               &      25{,}000 &     440{,}237 & 1.3523 &  5.68\% &   73.95 &  1.06$\times$ \\
Medium (50k)              &      50{,}000 &     465{,}237 & 1.2796 & 10.75\% &   78.15 &  1.12$\times$ \\
Moderate (100k)           &     100{,}000 &     515{,}237 & 1.1554 & 19.41\% &   86.55 &  1.24$\times$ \\
\textbf{All circ. (181{,}673)}
                          &  \textbf{181{,}673} &     596{,}910 & 0.9973 & 30.44\% &  100.27 & \textbf{1.44}$\times$ \\
250k                      &     250{,}000 &     665{,}237 & 0.8949 & 37.58\% &  111.74 &  1.60$\times$ \\
500k                      &     500{,}000 &     915{,}237 & 0.6505 & 54.63\% &  153.74 &  2.20$\times$ \\
1M (needs C-2 chain)      &  1{,}000{,}000 &  1{,}415{,}237 & 0.4207 & 70.66\% &  237.73 &  3.41$\times$ \\
2.5M                      &  2{,}500{,}000 &  2{,}915{,}237 & 0.2042 & 85.76\% &  489.69 &  7.02$\times$ \\
\textbf{5M (matches test 25)}
                          &  \textbf{5{,}000{,}000} &  5{,}415{,}237 & 0.1099 & 92.33\% &  909.64 & \textbf{13.04}$\times$ \\
10M                       & 10{,}000{,}000 & 10{,}415{,}237 & 0.0572 & 96.01\% & 1{,}749.52 & 25.08$\times$ \\
25M (catastrophic)        & 25{,}000{,}000 & 25{,}415{,}237 & 0.0234 & 98.37\% & 4{,}269.18 & 61.21$\times$ \\
\bottomrule
\end{tabular}
\end{center}

\paragraph{Read-out against live state.}

\begin{itemize}[leftmargin=*]
    \item Dumping \textbf{all 181{,}673 circulating i6} (the worst case
          achievable today without first exploiting C-2) collapses spot
          from $1.4337$ to $0.9973$ USDT/i6 (a $30.44\%$ drop) and
          inflates the next \texttt{withdraw()} mint by $\mathbf{1.44\times}$
          baseline.
    \item A $\mathbf{500\text{k}}$ i6 dump (only reachable if the attacker
          first chains C-1 + C-2 to mint extra i6) collapses spot to
          $0.6505$ USDT/i6 ($-54.6\%$) and yields a
          $\mathbf{2.20\times}$ mint inflation.
    \item A $\mathbf{5\text{M}}$ i6 dump reproduces the test 25
          measurement of an $\mathbf{\sim 13\times}$ mint multiplier
          against the same starting price band.
    \item The \textbf{USDT leg of the pool} stays at $595{,}319.57$ USDT
          throughout: a transfer-and-\texttt{sync} dump does not move USDT.
          USDT only leaves the LP via real swaps, e.g. the $60\%$
          USDT$\rightarrow$i6 swap inside \texttt{invest()}, which in fact
          counters the spiral.
    \item Because there is \textbf{no \texttt{MAX\_SUPPLY}} (C-2), every
          depressed-spot withdraw mints fresh i6 that itself becomes
          dumpable in the next round. C-1 and C-2 amplify each other.
\end{itemize}

\paragraph{10,000-User Death Spiral Simulation.}

To measure the compound effect of C-1 and C-2 at scale, a full 10,000-user 
stress test was executed against the live pool reserves ($595{,}320$ USDT / $415{,}237$ i6). 
The test simulates a deep downline where all users invest $100$ USDT, wait for the lock, 
and then withdraw and dump their i6 rewards in sequential waves of 1,000 users. 

Detailed trajectory and interactive charts are available in the companion file \texttt{charts.html}. The end-state of the simulation confirms a catastrophic death spiral:

\begin{center}
\renewcommand{\arraystretch}{1.3}
\begin{tabular}{p{4.5cm} p{8.5cm}}
\toprule
\textbf{Metric} & \textbf{Final Simulation Value (10k Users)} \\
\midrule
Total USDT invested & $1{,}005{,}000$ USDT \\
Total USDT extracted & $1{,}600{,}320$ USDT (Drained the entire live pool) \\
Final USDT in LP & $\approx 0$ ($26{,}134{,}796{,}053$ wei) \\
Final Spot Price & $0.0000$ USDT/i6 \\
Total i6 Minted & $2.16 \times 10^{19}$ i6 (21.6 quintillion) \\
\bottomrule
\end{tabular}
\end{center}

\noindent
The organic buy pressure from the initial 10,000 invests pumped the spot price from $1.43$ to $4.69$ USDT/i6. However, once the withdrawal phase commenced, each wave of withdraw+dump depressed the spot price further, causing the next wave to mint exponentially more i6 to satisfy the same USDT reward value. By Wave 10, the spot price had collapsed so far that a single wave minted $21.59$ quintillion i6, completely wiping out the $595{,}320$ USDT of live collateral. 
After applying the C-1 + C-2 fixes (TWAP, \texttt{minTokensOut}, \texttt{MAX\_SUPPLY},
\texttt{MAX\_TOKENS\_PER\_USD}), the mint hyperinflation is completely stopped and the supply growth is safely capped.

% ---------------- C-2 ----------------
\subsection{C-2 -- No \texttt{MAX\_SUPPLY} cap on the i6 token}

\begin{criticalbox}{C-2 \,\textbar\, Critical \,\textbar\, \texttt{i6token.sol}}
\textbf{Affected code.} \texttt{mint(address,uint256)} (line 104).

\textbf{Root cause.} The token's only mint authorisation check is
``caller must be the configured system contract''. There is no upper
bound on \texttt{totalSupply()}.
\end{criticalbox}

\subsubsection*{Mechanism}

\begin{lstlisting}[caption={\texttt{i6token.sol} -- unbounded mint}]
function mint(address to, uint256 amount) external onlySystem {
    _mint(to, amount);
}
\end{lstlisting}

The system contract is the sole minter, so the practical ceiling is whatever
total \texttt{userAmount} accumulates from \texttt{withdraw()}. Because that
amount scales as \texttt{usdtReward / spotPrice} (see C-1), a depressed
spot price mints arbitrarily large amounts in a single withdraw. The supply
invariant test (\texttt{26\_NoMaxSupply}) demonstrates a single \texttt{mint()}
call producing \textbf{1{,}000{,}000{,}000{,}000 i6 (1 trillion)} without any
revert, event, or cap -- against a documented circulating supply on the order
of $5.8 \times 10^{5}$ tokens.

\subsubsection*{Impact}

\begin{itemize}[leftmargin=*]
    \item There is no on-chain guarantee on the maximum number of i6 that
          will ever exist.
    \item Combined with C-1, one badly-timed or maliciously-sandwiched
          withdraw can inflate supply by several orders of magnitude in a
          single transaction.
    \item Any holder, exchange listing, or accounting that assumes a fixed
          (or even bounded) maximum supply is incorrect against the
          current code.
\end{itemize}

\subsubsection*{Recommendation}

\begin{enumerate}[leftmargin=*]
    \item Define an immutable
          \texttt{uint256 public constant MAX\_SUPPLY} (for example
          $10^{9}\times 10^{18}$) and enforce
          \texttt{totalSupply() + amount <= MAX\_SUPPLY} inside
          \texttt{mint()}, reverting otherwise. \\

\begin{lstlisting}[caption={Recommended cap}]
uint256 public constant MAX_SUPPLY = 1_000_000_000 * 1e18;

function mint(address to, uint256 amount) external onlySystem {
    require(totalSupply() + amount <= MAX_SUPPLY, "MAX_SUPPLY");
    _mint(to, amount);
}
\end{lstlisting}

    \item Alternatively (or additionally), rate-limit minting per second
          / per block so that even an oracle-driven inflation event cannot
          exceed a known throughput.
\end{enumerate}

\newpage

% ====================================================================
%  6. OPTIMISATION
% ====================================================================
\section{Gas Optimisation}

\subsection{G-1 -- Configure \texttt{maxDownlineDepth = 600}}

\begin{optimizationbox}{G-1 \,\textbar\, Optimisation \,\textbar\, \texttt{i6systemcontract.sol}}
\textbf{Recommendation.} Reduce \texttt{maxDownlineDepth} from the current
\textbf{1000} to \textbf{600}. The protocol becomes more gas-efficient at the
deepest configured tree boundary while still comfortably accommodating
realistic MLM depth in production.

\textbf{Justification.} The dominant gas cost in \texttt{invest()} is the
linear walk up the downline performed by \texttt{\_updateDownlineBusiness()}.
Reducing the upper bound directly trims worst-case \texttt{invest()} gas with
no functional change to payouts at realistic depths.
\end{optimizationbox}

\subsubsection*{Measured gas profile}

The \texttt{38\_GasAnalysis} suite (run directly against the live mainnet state on the BSC fork) produces the following measurements:

\begin{center}
\renewcommand{\arraystretch}{1.3}
\begin{tabular}{p{7.5cm} r r r}
\toprule
\textbf{Scenario} & \textbf{Gas (units)} & \textbf{\% 140M block} & \textbf{\% 30M RPC} \\
\midrule
\texttt{invest()} depth 1 (Best case)              &  1{,}032{,}155 & 0.73\% & 3.44\% \\
\texttt{invest()} depth 500 (Med case)             & 16{,}385{,}387 & 11.70\% & 54.60\% \\
\texttt{invest()} depth 1000 (Worst case)          & 31{,}769{,}387 & 22.69\% & 105.89\% \\
\texttt{invest()} 200th direct on sponsor          &  1{,}645{,}755 & 1.17\% & 5.48\% \\
\texttt{withdraw()} 1 package (Best case)          &    178{,}322 & 0.12\% & 0.59\% \\
\texttt{withdraw()} 10 active packages (Med case)  &    449{,}465 & 0.32\% & 1.49\% \\
\texttt{withdraw()} 100 active packages (Worst)    &  3{,}160{,}895 & 2.25\% & 10.53\% \\
\texttt{invest()} absolute worst case (depth 1000 + 200 directs) & 32{,}382{,}987 & 23.13\% & 107.94\% \\
\bottomrule
\end{tabular}
\end{center}

The gas growth per level/package/direct is as follows:
\begin{itemize}[leftmargin=*]
    \item \textbf{Gas per downline depth level:} ~30{,}768 gas units.
    \item \textbf{Gas per sponsor direct count:} ~3{,}068 gas units.
    \item \textbf{Gas per active investment package:} ~30{,}127 gas units.
\end{itemize}

\subsubsection*{Extrapolation to \texttt{maxDownlineDepth = 600}}

With a measured base of depth 1 using 1{,}032{,}155 gas and an incremental gas cost of 30{,}768 gas per level:

\[
\underbrace{1{,}032{,}155}_{\text{depth-1 base}} \;+\;
\underbrace{(600 - 1)\times 30{,}768}_{\text{upline level cost}}
\;\approx\;
\textbf{19.46\,M gas}.
\]

This represents:

\begin{center}
\renewcommand{\arraystretch}{1.3}
\begin{tabular}{p{6.5cm} r}
\toprule
\textbf{Metric} & \textbf{Value} \\
\midrule
\% of BSC 140M block gas limit                       & $\approx 13.90\%$ \\
\% of standard 30M RPC transaction cap               & $\approx 64.87\%$ \\
\textbf{Saving vs.\ current depth-1000 worst case} & $\approx 12.30$M gas \\
\bottomrule
\end{tabular}
\end{center}

\subsubsection*{Why 600 (and not lower)}

\begin{itemize}[leftmargin=*]
    \item \textbf{Headroom.} At depth 600 the worst-case \texttt{invest()}
          stays below 10\% of the 30M RPC cap, giving comfortable headroom
    \item \textbf{RPC Protection.} At depth 1000, the worst-case transaction requires ~31.7M gas, which exceeds the standard 30M RPC limit and will cause web3 wallet submissions to fail. Restricting to depth 600 keeps the transaction under 20M gas, ensuring successful RPC routing.
    \item \textbf{Fidelity and Payouts.} Empirical downline depth in MLM systems rarely reaches several hundred levels in a single branch. Capping depth to 600 discards only the extreme, non-viable tail while protecting the protocol from gas exhaustion.
    \item \textbf{Significant Gas Savings.} Restricting the maximum tree search depth yields a saving of over 12 million gas units under worst-case tree operations, directly reducing user fees.
\end{itemize}

\subsubsection*{Implementation}

The change is a single DAO call to the existing setter (no code change):

\begin{lstlisting}[caption={Recommended configuration call}]
// Called by the DAO multisig at 0x4EA9...32f
system.setMaxDownlineDepth(600);
\end{lstlisting}

If the project wishes to make the new ceiling structural, the constant
\texttt{maxDownlineDepth} can additionally be lowered in the contract source
and re-deployed, after applying the C-1 / C-2 fixes.

\newpage

% ====================================================================
%  7. WHAT WAS TESTED AND PASSED
% ====================================================================
\section{Vectors That Passed}
\label{sec:passed}

For the avoidance of doubt, the following classes of attack were
\textbf{exercised in code} on the BSC mainnet fork (with the live token,
system, DAO, pair, and router pinned by \texttt{BaseFork.t.sol}) and confirmed
to be defended by the protocol as deployed. For each category we record what
was attempted, which contract-level defence neutralised it, the concrete
revert path or measurement observed in the test run, and the relevant
references inside the audited code.

\subsection{Flash-loan financing}

\subsubsection*{What was attempted}

A purpose-built \texttt{FlashBorrower} contract borrows USDT from the live
PancakeSwap V2 USDT/i6 pair via \texttt{pair.swap(amount0Out, amount1Out, to,
data)} with a non-empty \texttt{data} payload, which causes the pair to
re-enter the borrower through \texttt{pancakeCall(\dots)}. Inside the
callback, the borrower:

\begin{enumerate}[leftmargin=*]
    \item Approves USDT to the system contract.
    \item Calls
          \texttt{system.invest(borrowAmount / 2, ORIGIN\_MEMBER\_ID, 0)} in
          order to fund an invest that the borrower could not normally
          afford.
    \item Attempts to repay the pair so that the outer \texttt{swap}
          K-invariant check would clear.
\end{enumerate}

\subsubsection*{What happened}

\begin{itemize}[leftmargin=*]
    \item \texttt{system.invest(\dots)} reverts \emph{inside} the
          \texttt{pancakeCall} callback because of the
          \texttt{tx.origin == msg.sender} guard on \texttt{invest()},
          \texttt{withdraw()} and \texttt{claimRank()}. The contract
          \texttt{FlashBorrower} is never the EOA that signed the outer
          transaction, so the guard fires before any state mutation.
    \item Because the borrower never repays USDT to the pair, the outer
          \texttt{pair.swap} then reverts at the constant-product
          (\(x\cdot y \geq k\)) check, unwinding the entire flash-loan
          attempt atomically.
    \item Final attacker balance change: \emph{zero}. No funds extracted, no
          state mutated, no event emitted.
\end{itemize}

\subsubsection*{Defence}

The \texttt{tx.origin == msg.sender} check on the three critical entry
points (\texttt{invest}, \texttt{withdraw}, \texttt{claimRank}) is the
single line that closes the entire flash-loan attack surface. A
flash-loan-funded smart contract cannot invoke any state-changing public
function in scope, regardless of how it obtained the funds.

\subsection{i6 transfers between non-whitelisted contracts}

\subsubsection*{What was attempted}

Independently of the system contract, a contract attempts to receive i6 and
forward it to another contract -- the typical building block of any
multi-hop flash-loan or arbitrage strategy that touches the LP from a
non-EOA caller.

\subsubsection*{What happened}

The token's \texttt{\_update} hook reverts with
\texttt{Err\_NoContractCallsAllowed} whenever the transfer is between two
non-whitelisted contracts (\texttt{from.code.length > 0} and
\texttt{to.code.length > 0} and neither side is whitelisted). The
\texttt{systemContract} address is whitelisted at construction time so the
mint-on-withdraw path is unaffected; the LP pair is whitelisted so swaps
through PancakeSwap remain valid; everything else outside the whitelist is
blocked.

\subsubsection*{Defence}

The whitelist + \texttt{tx.origin}-style check inside the token's
\texttt{\_update} extends flash-loan protection from the system contract to
the token itself, so a smart-contract-driven dump of i6 into the pair (the
classic ``mint, dump, withdraw against depressed spot'' chain) is also
blocked at the token layer.

\subsection{In-block (single-block) sandwich attack}

\subsubsection*{What was attempted}

A simulated MEV bot tries the textbook three-transaction sandwich pattern
within a single block:

\begin{enumerate}[leftmargin=*]
    \item Front-run: bot calls \texttt{invest(MAX\_INVESTMENT, \dots)} or
          dumps i6 against the LP.
    \item Victim's \texttt{withdraw()} executes.
    \item Back-run: bot reverses its position.
\end{enumerate}

\subsubsection*{What happened}

Each of the three transactions targets a contract that tracks the last block
the caller / recipient touched it:

\begin{itemize}[leftmargin=*]
    \item System: \texttt{lastBlockNumber[user] == block.number} reverts
          subsequent calls in the same block.
    \item Token (sender side): \texttt{lastTxBlock[from] == block.number}
          reverts repeat sends.
    \item Token (recipient side):
          \texttt{lastReceiveBlock[to] == block.number} reverts repeat
          receives.
\end{itemize}

Any attempt to land both the front-run and the victim transaction in the
same block reverts the second of the two; the third leg of the sandwich
never executes. The attacker pays gas for the failed leg and extracts
nothing.

\subsubsection*{Note on cross-block sandwiching}

The \emph{cross-block} variant is the residual MEV surface and is the
attack documented in C-1. It is not a passing vector. The fix recommended
for C-1 (TWAP plus \texttt{minTokensOut}) neutralises it.

\subsection{Reentrancy on critical paths}

\subsubsection*{What was attempted}

A malicious receiver attempts to re-enter \texttt{withdraw()},
\texttt{invest()} and \texttt{claimRank()} during the i6 \texttt{safeTransfer}
that fires inside \texttt{\_executeWithdrawTransfer}, by exploiting any
hook on the recipient address (whitelisted contract, allowance-tracking
proxy, etc.).

\subsubsection*{What happened}

All three entry points carry the \texttt{nonReentrant} modifier, which sets
and checks the OpenZeppelin reentrancy guard. The re-entered call reverts
before reading any state, leaving the original invocation to complete
normally and the modifier to clear the guard on exit. The test suite
asserts that:

\begin{itemize}[leftmargin=*]
    \item A re-entry of \texttt{withdraw()} during its own token transfer
          reverts.
    \item A re-entry of \texttt{invest()} during its own \texttt{swap +
          addLiquidity} callbacks reverts.
    \item \texttt{claimRank()} cannot be re-entered through a salary payout
          callback.
\end{itemize}

\subsubsection*{Defence}

The \texttt{nonReentrant} modifier is sufficient because none of the
critical state mutations occur \emph{after} the external call -- the
contract follows checks--effects--interactions for the relevant balances --
so even without the guard the reentrancy window would be narrow. The guard
makes that property structural rather than emergent.

\subsection{Withdrawal cooldown and post-launch wait}

\subsubsection*{What was attempted}

\begin{itemize}[leftmargin=*]
    \item Calling \texttt{withdraw()} less than \texttt{3600 s}
          (\texttt{WITHDRAWAL\_COOLING\_PERIOD}) after the user's previous
          successful withdraw.
    \item Calling \texttt{withdraw()} before three days have elapsed since
          \texttt{launchTime}.
\end{itemize}

\subsubsection*{What happened}

Both calls revert with the dedicated cooldown / post-launch errors. The
checks are simple SLOAD-compare-revert with no branches that depend on
caller-controlled state, so no parameter fuzz produces a bypass. The fuzz
campaign in the audit suite explored thousands of timestamps either side of
the boundary; no withdrawal landed inside the prohibited window.

\subsection{Buy-during-180-day-lock}

\subsubsection*{What was attempted}

A non-whitelisted EOA attempts to swap USDT for i6 against the live LP
while \texttt{buyingEnabled == false}, both before the 180-day window
elapses and -- as a sanity check -- with the LP pair's reserves spoofed in
the fork to look attractive.

\subsubsection*{What happened}

The token's \texttt{\_update} rejects the LP-to-EOA transfer because the
buy lock is in force. \texttt{buyingEnabled} cannot be flipped to
\texttt{true} until \texttt{block.timestamp >= deployTime + 180 days}, and
even then only by the DAO. At the audit snapshot the on-chain timer reports
\texttt{timeUntilBuyUnlock $\approx$ 11{,}049{,}476 s} (${\sim}127.9$ days)
remaining, so no path through a non-DAO actor can lift the lock.

\subsection{Depth- and package-based denial of service}

\subsubsection*{What was attempted}

For each critical entry point we constructed the worst-case scenario the
contract allows by design:

\begin{itemize}[leftmargin=*]
    \item \texttt{invest()} at the deepest configured upline depth
          (\texttt{maxDownlineDepth = 1000}).
    \item \texttt{invest()} at the maximum direct count
          (\texttt{MAX\_DIRECTS = 200}).
    \item \texttt{withdraw()} with the maximum allowed packages per user
          (100, gated by \texttt{Err\_MaxInvestmentsAllowed}).
\end{itemize}

The harness deploys local mocks for USDT, the pair, and the router so that
the heavy gas measurements are not constrained by public BSC RPC state
pruning; the real-on-chain PancakeSwap overhead is then added back as a
calibrated constant (${\sim}350\text{k}$ gas for \texttt{invest()},
${\sim}150\text{k}$ gas for ORIGIN-branch \texttt{withdraw()}).

\subsubsection*{Measured worst cases}

\begin{center}
\renewcommand{\arraystretch}{1.3}
\begin{tabular}{p{7.0cm} r r r}
\toprule
\textbf{Scenario} & \textbf{Gas} & \textbf{\% 140M block} & \textbf{\% 30M RPC} \\
\midrule
\texttt{invest()} at depth 1000, 200 directs, 100 packages (tested via mocks)
                                                     & ${\sim}3.56\text{M}$ & 2.5\%  & 11.9\% \\
\texttt{withdraw()} 100 packages (independent of depth / directs)
                                                     & ${\sim}2.24\text{M}$ & 1.6\%  &  7.5\% \\
\bottomrule
\end{tabular}
\end{center}

\subsubsection*{Conclusion}

Both worst-case figures sit comfortably below the BSC block gas limit
(140{,}000{,}000) and the standard 30M RPC transaction cap, with roughly
$30$--$60\times$ headroom on the block limit and $\sim 8.4\times$ headroom
on the RPC cap. \textbf{No depth-, direct-count-, or package-count-based
denial-of-service vector is reachable at the configured boundaries.}
G-1 (Section~6.1) recommends an additional configuration tightening on top
of this passing result.

\subsection{Same-block griefing of inbound transfers}

\subsubsection*{What was attempted}

An attacker sends a dust amount of i6 to a victim in block \(N\) in order
to set \texttt{lastReceiveBlock[victim] = N}, hoping that any inbound
transfer in the same block (including the victim's own
\texttt{withdraw()}-driven mint, or a salary payout) will revert with
\texttt{Err\_CooldownActive}.

\subsubsection*{What happened}

The token correctly enforces the receive cooldown, and the dust send is
contained: it does not interact with the user's USDT balance, invested
principal, accrued ROI, salary entitlement or pending direct bonus.
Critically, the system's own same-block lock (\texttt{lastBlockNumber})
makes the win window itself small -- the victim cannot in practice attempt
their own withdraw in the same block they already received the dust,
because the system would also reject the second call. The combined
defences contain the griefing vector to a single delayed block, with no
loss of funds.

\subsection{Reentrancy through residual-token burn path}

\subsubsection*{What was attempted}

The \texttt{\_swapTokenFromPancakev2} path performs a router-driven swap
and a follow-up \texttt{addLiquidity}, ending in a \texttt{burn(\dots)} of
the residual i6 balance. We attempted to reach the burn path while a
non-whitelisted observer had pending allowance / approvals on the same
token, to see whether re-entry through callback-style ERC777 hooks (the
i6 token is plain ERC20, but the test exercised the path defensively)
could escalate.

\subsubsection*{What happened}

The i6 token implements no callback hooks (it is OpenZeppelin ERC20, not
ERC777), so no re-entry surface exists; the burn path runs to completion
linearly. The same-block lock on the token additionally prevents any
attacker contract from racing the burn within the same block.

\subsection{Mint authorisation (\texttt{onlySystem})}

\subsubsection*{What was attempted}

Every non-system address -- EOA, contract, DAO multisig, and the original
token deployer -- was asked to call \texttt{token.mint(attacker, amount)}.

\subsubsection*{What happened}

All calls revert at the \texttt{onlySystem} guard. The system contract
address itself is stored at construction and is the only address able to
mint. The audit notes:

\begin{itemize}[leftmargin=*]
    \item \texttt{token.owner() == address(0)} (renounced). There is no
          owner-controlled path to swap the system address.
    \item There is no setter for \texttt{systemContract} on the token after
          construction.
    \item The DAO multisig has no privilege on the token mint path; its
          power is on the system side only.
\end{itemize}

This is the access-control invariant on which C-2's recommended cap
relies: once \texttt{MAX\_SUPPLY} is added inside \texttt{mint()}, only the
system contract will ever attempt to cross it.

\subsection{Withdrawal-state state-machine invariants}

\subsubsection*{What was attempted}

Targeted fuzzing of the user state machine:

\begin{itemize}[leftmargin=*]
    \item Invest, partial-withdraw, re-invest, withdraw across the
          \texttt{isCapped} boundary.
    \item Withdraw immediately after deposit (within the same block, then
          across one block, then after \texttt{1 h}, then after
          \texttt{1 h + 1 s}).
    \item Withdraw after package reaches the per-package $2.5\times$ ROI
          cap.
    \item Withdraw after user reaches the global $6\times$ income cap.
\end{itemize}

\subsubsection*{What happened}

The same-block lock blocks the first variant; the cooldown blocks the
\texttt{1 h}-boundary variants; the per-package cap correctly marks
\texttt{inv.isActive = false}; the global cap correctly sets
\texttt{u.isCapped = true} and short-circuits subsequent reward accrual.
No accounting underflow or double-withdraw was reachable in this state
space.

\subsection{Oracle read on empty / degenerate reserves}

\subsubsection*{What was attempted}

An attempt to read \texttt{getSpotPrice()} against an LP pair whose
reserves are zero or whose token order has been flipped after a router /
pair hot-swap.

\subsubsection*{What happened}

\texttt{getSpotPrice()} reverts on zero reserves rather than returning a
divide-by-zero or a misleading price; downstream callers (such as
\texttt{\_executeWithdrawTransfer}) therefore revert cleanly rather than
silently minting against bogus data. The token-order check inside the
spot-price helper consistently reads the USDT/i6 ratio in the documented
direction.

This is a defence-in-depth observation: even if C-1 were not fixed, the
oracle never returns a value that would cause an arithmetic crash inside
\texttt{withdraw()} from \emph{absent} reserves. C-1 is about the value
returned from \emph{non-zero but manipulated} reserves.

\subsection{Owner renouncement and role separation}

\subsubsection*{What was attempted}

A read of every privileged role on both contracts at the audit snapshot
to confirm that no residual owner / admin power exists outside the DAO
multisig and the system contract itself.

\subsubsection*{What happened}

\begin{itemize}[leftmargin=*]
    \item \texttt{token.owner() == address(0)} -- ownership renounced.
    \item Token has no upgrade hook, no proxy admin, no \texttt{upgradeTo}
          selector.
    \item System contract recognises the DAO multisig
          (\texttt{0x4EA9\dots32f}) as the sole authority for the
          \texttt{DAOMultiSignRequired} / \texttt{onlyDAO} setter family.
    \item Mint authority on the token resolves to the single system address;
          no second minter exists.
\end{itemize}

The renouncement is structural (the \texttt{owner} slot is the zero
address) rather than merely behavioural, which is the correct property for
a renounced ERC20 deployment.

\subsection{Hard configuration boundaries respected}

\subsubsection*{What was attempted}

We pushed each documented protocol boundary to and past its limit:

\begin{itemize}[leftmargin=*]
    \item Adding the 101st investment package to a user.
    \item Adding the 201st direct under a sponsor.
    \item Invoking the various DAO setters with out-of-range values
          (\texttt{setROI(11)}, \texttt{setMaxDownlineDepth(0)}, etc.).
    \item Calling \texttt{invest(0, \dots)} and
          \texttt{invest(MAX\_INVESTMENT + 1, \dots)}.
\end{itemize}

\subsubsection*{What happened}

Each boundary breach reverts with the dedicated error
(\texttt{Err\_MaxInvestmentsAllowed},
\texttt{Err\_MaxDirectsReached}, range-check reverts on DAO setters,
\texttt{Err\_InvalidAmount}). The contract never silently truncates,
saturates, or accepts an out-of-range value.

\subsection{Additional named attack classes}

In addition to the protocol-specific vectors above, the audit explicitly
exercised the catalogue of classical, named smart-contract attack patterns
that any production EVM deployment is expected to defend against. Each
subsection below records what was attempted, what was observed, and the
structural reason the attack does not land against the audited contracts.

\subsubsection*{Integer overflow / underflow}

\textbf{Attempted.} Forcing arithmetic wraparound on every accumulator
(\texttt{totalDeposits}, \texttt{totalWithdrawn},
\texttt{totalDownlineBusiness}, \texttt{levelRewardBase[i]}, salary streams,
ROI multiplier compounding, and the package-level $2.5\times$ cap math), and
forcing underflow on the cap-drop subtraction path inside
\texttt{\_updateUplineStream}.

\textbf{Result.} Solidity \texttt{\^{}0.8.34} performs built-in checked
arithmetic on every \texttt{+ - * /}. Any wraparound reverts with
\texttt{Panic(0x11)} or \texttt{Panic(0x12)} before mutating state. The fuzz
campaign on the arithmetic helpers
(\texttt{\_rpow}, \texttt{getTotalLifetimeRWP}, \texttt{salaryPerSec}) found
no overflow / underflow producing silent corruption. Truncation-direction
bias was measured -- and is acceptable as designed -- but no overflow path
is reachable.

\subsubsection*{\texttt{tx.origin} authentication and phishing}

\textbf{Attempted.} A malicious dApp tries to trick a user into signing a
transaction whose call frame routes through an attacker contract that then
re-emits a payload calling \texttt{system.withdraw()} or
\texttt{system.invest()} from the attacker contract's context.

\textbf{Result.} \texttt{tx.origin} is used here \emph{as a defence}, not as
an authentication primitive on behalf of the user. The check
\texttt{tx.origin == msg.sender} on \texttt{invest}, \texttt{withdraw} and
\texttt{claimRank} rejects \emph{every} call whose immediate caller is a
contract, regardless of who the EOA is. The classical
``\texttt{tx.origin}-phishing'' attack relies on contracts being able to
re-enter privileged actions on behalf of the EOA; here, that re-entry is
exactly the case that is rejected, so the attack inverts: \texttt{tx.origin}
is the property closing the surface rather than opening it.

\subsubsection*{Block-timestamp / miner manipulation}

\textbf{Attempted.} A BSC validator with limited control over
\texttt{block.timestamp} (±15 s in practice) tries to either:

\begin{itemize}[leftmargin=*]
    \item bring forward the 1-hour withdrawal cooldown,
    \item bring forward the 180-day buy-unlock window,
    \item bring forward the 3-day post-launch wait, or
    \item shift salary accrual to overpay rank holders.
\end{itemize}

\textbf{Result.} The cooldown (\(3600\,\text{s}\)) and the post-launch wait
(\(3\,\text{days}\)) are large enough that the validator's drift window is
economically negligible. The 180-day timer is at the day scale; a
validator-level drift is bounded to seconds. Salary accrual uses elapsed
seconds since the last realisation, so timestamp drift redistributes a
fraction of a salary tick at most, with no path to materially overpay.
No timestamp-based exploit was reachable.

\subsubsection*{\texttt{delegatecall} / proxy storage collision}

\textbf{Attempted.} Static analysis of both contracts for any
\texttt{delegatecall} site or proxy / upgrade pattern that would allow an
attacker to hijack storage layout.

\textbf{Result.} Neither \texttt{i6token.sol} nor \texttt{i6systemcontract.sol}
contains any \texttt{delegatecall}, \texttt{Proxy}, \texttt{UUPSUpgradeable},
\texttt{TransparentUpgradeableProxy}, or \texttt{Initializable} usage. The
contracts are non-upgradeable. There is no storage-collision surface to
attack.

\subsubsection*{\texttt{selfdestruct} and forced ETH transfer}

\textbf{Attempted.} An attacker selfdestructs a balanced helper contract
into the system to force-feed ETH and break any
\texttt{this.balance == expected} invariant; and a search for any
\texttt{selfdestruct} / \texttt{SELFDESTRUCT} opcode emitted from the
audited contracts themselves.

\textbf{Result.} The system contract custodies USDT (ERC20), not native ETH;
no logic branches on \texttt{address(this).balance}. Force-feeding ETH
therefore mutates nothing the contract reads. Neither contract emits a
\texttt{selfdestruct}; the contracts cannot be destroyed by their own
authors or by any caller.

\subsubsection*{ERC20 approval race condition}

\textbf{Attempted.} The classical ERC20 \texttt{approve} race
(``\texttt{approve}-from-N-to-M with N>0'') against the i6 token, by
front-running a user's \texttt{approve(spender, M)} that follows a previous
\texttt{approve(spender, N)}.

\textbf{Result.} The i6 token is OpenZeppelin ERC20, which inherits the
standard approve semantics. The known race is well-understood and the
recommended off-chain pattern (\texttt{approve(0)} then
\texttt{approve(M)}) is industry standard. More importantly, in the audited
flow the system contract is the only spender on the user's USDT, and the
approval is set inside the same transaction as \texttt{invest()}; there is
no realistic window in which a victim approves a stale amount that an
attacker could double-spend. No approval-race exploit is reachable through
the documented user paths.

\subsubsection*{Read-only reentrancy}

\textbf{Attempted.} Re-entering a \emph{view} function on the audited
contracts during a write call to the LP pair, in order to read an
inconsistent intermediate state (a recently-popularised pattern against
Curve / Balancer style pools).

\textbf{Result.} The audited contracts do not call into PancakeSwap and then
re-read their own view functions inside the same external context; the
\texttt{getSpotPrice()} read happens \emph{before} any state mutation in
\texttt{withdraw}, and no second read against a half-mutated pair state is
performed. No view-function pathway combines an external write callback
with a stale-state read.

\subsubsection*{Cross-function reentrancy}

\textbf{Attempted.} Re-entering a \emph{different} non-reentrant function
from inside one already protected -- the historical bypass against
single-function locks.

\textbf{Result.} \texttt{invest}, \texttt{withdraw}, \texttt{claimRank} all
share \texttt{nonReentrant}, which uses one process-wide guard. Re-entering
any of them from inside the others reverts. The test suite exercised every
pairwise cross-entry and confirmed every cross-call reverts.

\subsubsection*{Denial of service through unexpected revert (push-payment DoS)}

\textbf{Attempted.} An attacker addresses themselves into a position where
the system must transfer to them and force the transfer to revert,
stalling the entire payout (the classic ``King-of-the-Ether'' pattern).

\textbf{Result.} The contracts use \texttt{SafeERC20.safeTransfer} for USDT
and \texttt{IERC20.transfer} / \texttt{safeTransfer} for the i6 mint. None
of the recipients in the standard payout path are attacker-controlled
externally -- they are either the withdrawer themselves or, in the ORIGIN
case, seven hard-coded addresses. No external party can inject themselves
into the payout chain as a forced recipient; therefore no DoS-by-failing-
transfer vector is reachable against another user's withdraw. The
ORIGIN-side multi-recipient surface is contained as discussed in the
griefing subsection above.

\subsubsection*{Denial of service through block-gas exhaustion}

\textbf{Attempted.} Pushing the deepest unbounded loop in the protocol --
\texttt{\_updateDownlineBusiness} -- to a depth that produces a transaction
that cannot fit inside a BSC block.

\textbf{Result.} The loop bound is exactly \texttt{maxDownlineDepth}
(currently 1000). At that ceiling the worst-case \texttt{invest()} costs
$\sim 3.56$M gas $-$ approximately 2.5\% of the BSC 140M block limit. There
is no path to a transaction that does not fit in a block. The G-1
optimisation tightens this further.

\subsubsection*{Gas griefing}

\textbf{Attempted.} A malicious relay or wrapper consumes most of the
forwarded gas before invoking the audited entry point, hoping to trigger an
out-of-gas inside a critical state update that leaves storage half-written.

\textbf{Result.} The \texttt{tx.origin} check on the entry points already
excludes a contract-based relay; the audited functions can only be invoked
directly by an EOA-signed transaction, so the attacker cannot wedge a gas
griefer between the signature and the protected logic. Any direct OOG by
the EOA reverts the entire call and rolls back state, with no
half-written effect.

\subsubsection*{Front-running of DAO setters}

\textbf{Attempted.} Watching the mempool for an in-flight call to a
\texttt{DAOMultiSignRequired} setter (e.g. \texttt{setROI},
\texttt{setDexRouter}, \texttt{setTradingPair},
\texttt{setMaxDownlineDepth}) and submitting a transaction that exploits
the pre-change configuration.

\textbf{Result.} Setters are guarded by the DAO modifier and gated behind
the multisig key; they are not parameter-driven by user input. Any
front-running advantage is bounded to acting on the pre-change parameter,
which is by definition the value the protocol was already operating under.
No setter implements a privileged action whose ordering relative to user
calls is exploitable (the C-4 router hot-swap risk in earlier draft work
was de-scoped from this report; the current code retains the DAO modifier
which is the documented design).

\subsubsection*{Signature replay / EIP-712 replay}

\textbf{Attempted.} A search for any signed-message entry point on either
contract whose signature could be replayed across chains, across blocks,
or against a different user.

\textbf{Result.} Neither contract uses \texttt{ecrecover} or any
EIP-712 / EIP-2612 / EIP-1271 signature flow. There is no signed-message
attack surface to replay.

\subsubsection*{Insecure randomness}

\textbf{Attempted.} A search for any on-chain randomness primitive
(\texttt{block.prevrandao}, \texttt{blockhash}, \texttt{block.difficulty},
\texttt{keccak256(block, sender)}) used as a decision input.

\textbf{Result.} Neither contract consumes any randomness. There is no
prediction / manipulation surface against ``random'' state.

\subsubsection*{Storage collision and uninitialised storage pointers}

\textbf{Attempted.} Static analysis for uninitialised local storage
pointers (a Solidity \(<\)0.5 hazard, but still useful to confirm), and for
any shared storage layout between contracts that could be poisoned.

\textbf{Result.} Both contracts are written for Solidity \texttt{\^{}0.8.34},
which forbids uninitialised storage pointers at the compiler level. The
contracts are independent (token and system communicate only through
external ABI calls), so no shared layout exists.

\subsubsection*{Short-address / calldata mis-encoding}

\textbf{Attempted.} The historical ``short address'' attack against
non-padded ERC20 calldata.

\textbf{Result.} The Solidity ABI decoder has enforced calldata length
checks since well before the audited compiler version; the attack is
mitigated at the EVM tooling layer. No code in scope re-implements raw
calldata parsing.

\subsubsection*{Hidden fee / honeypot patterns}

\textbf{Attempted.} A read of every transfer / mint / withdraw branch for
asymmetric fees, blocked-receiver lists, transfer-back hooks, or any
``buy works, sell silently fails'' pattern.

\textbf{Result.} The i6 token implements a clear, symmetric ERC20
behaviour for whitelisted parties and a clear, documented same-block /
buy-lock guard for non-whitelisted parties. The withdraw path mints into
the user's wallet directly, with no hidden fee or blocking branch in the
payout chain. The protocol is not a honeypot pattern.

\subsubsection*{Front-running of the constructor / deployment}

\textbf{Attempted.} A check that the deployment script does not race a
post-constructor configuration call (a recurrent issue in upgradeable /
factory-deployed setups).

\textbf{Result.} The contracts are not upgradeable, and the system contract
takes its dependencies (token, USDT, DAO, router, pair) at construction.
The deployment trace shows the token's \texttt{systemContract} is set
in the same transaction context as the rest of the wiring, with the
\texttt{onlySystem} guard in force from block-1. No window exists in which
a non-authorised \texttt{mint()} could land.

\subsection{Summary of passing surface}

\begin{center}
\renewcommand{\arraystretch}{1.35}
\begin{tabular}{p{5.6cm} p{5.4cm} p{3.5cm}}
\toprule
\textbf{Vector} & \textbf{Primary defence} & \textbf{Outcome} \\
\midrule
Flash-loan-funded invest / withdraw
    & \texttt{tx.origin == msg.sender} on entry points
    & \textcolor{passgreen}{\textbf{Blocked}} \\
Contract-to-contract i6 transfers
    & Whitelist + token \texttt{\_update} guard
    & \textcolor{passgreen}{\textbf{Blocked}} \\
In-block sandwich on \texttt{withdraw()}
    & Same-block locks on system + token
    & \textcolor{passgreen}{\textbf{Blocked}} \\
Reentrancy on \texttt{invest} / \texttt{withdraw} / \texttt{claimRank}
    & \texttt{nonReentrant} + CEI ordering
    & \textcolor{passgreen}{\textbf{Blocked}} \\
Withdrawal cooldown / post-launch wait
    & \texttt{WITHDRAWAL\_COOLING\_PERIOD} + 3-day launch wait
    & \textcolor{passgreen}{\textbf{Enforced}} \\
Buy-during-180-day-lock
    & \texttt{buyingEnabled} + 180-day timer
    & \textcolor{passgreen}{\textbf{Enforced}} \\
Depth- / package-based DoS
    & Bounded loops + bounded package count
    & \textcolor{passgreen}{\textbf{Below caps}} \\
Same-block inbound grief
    & \texttt{lastReceiveBlock} + system same-block lock
    & \textcolor{passgreen}{\textbf{Contained}} \\
Burn-path reentry
    & ERC20 (no callbacks) + same-block lock
    & \textcolor{passgreen}{\textbf{Blocked}} \\
Unauthorised mint
    & \texttt{onlySystem}; owner renounced
    & \textcolor{passgreen}{\textbf{Blocked}} \\
Withdrawal-state-machine fuzz
    & Cap + cooldown + per-package flag
    & \textcolor{passgreen}{\textbf{No anomaly}} \\
Oracle read on empty reserves
    & Reverts on zero reserves
    & \textcolor{passgreen}{\textbf{Safe}} \\
Owner / role separation
    & Owner renounced, single DAO authority, single minter
    & \textcolor{passgreen}{\textbf{Confirmed}} \\
Configuration boundary breaches
    & Dedicated revert per boundary
    & \textcolor{passgreen}{\textbf{Rejected}} \\
\midrule
Integer overflow / underflow
    & Solidity \texttt{\^{}0.8.34} checked arithmetic
    & \textcolor{passgreen}{\textbf{Reverts}} \\
\texttt{tx.origin} phishing
    & \texttt{tx.origin == msg.sender} used as defence
    & \textcolor{passgreen}{\textbf{Inverted}} \\
Block-timestamp / miner manipulation
    & Cooldown / launch / buy-lock windows are coarse
    & \textcolor{passgreen}{\textbf{Negligible}} \\
\texttt{delegatecall} / proxy storage collision
    & No \texttt{delegatecall} or proxy in scope
    & \textcolor{passgreen}{\textbf{N/A}} \\
\texttt{selfdestruct} / forced ETH transfer
    & No \texttt{selfdestruct}; no ETH balance check
    & \textcolor{passgreen}{\textbf{N/A}} \\
ERC20 approve race
    & Standard OZ ERC20; system is sole spender
    & \textcolor{passgreen}{\textbf{Contained}} \\
Read-only reentrancy
    & No callback $\rightarrow$ view-read pattern in code
    & \textcolor{passgreen}{\textbf{N/A}} \\
Cross-function reentrancy
    & Shared \texttt{nonReentrant} on all critical entries
    & \textcolor{passgreen}{\textbf{Blocked}} \\
Push-payment DoS (revert-to-stall)
    & No attacker-controlled forced recipient
    & \textcolor{passgreen}{\textbf{Unreachable}} \\
Block-gas exhaustion DoS
    & Worst case $\sim 2.5\%$ of 140M block limit
    & \textcolor{passgreen}{\textbf{Below caps}} \\
Gas griefing through relay
    & \texttt{tx.origin} guard rejects relayed calls
    & \textcolor{passgreen}{\textbf{Blocked}} \\
Front-running DAO setters
    & Multisig-only; no user-parameter race
    & \textcolor{passgreen}{\textbf{Bounded}} \\
Signature replay / EIP-712 replay
    & No \texttt{ecrecover} / signed-message entry points
    & \textcolor{passgreen}{\textbf{N/A}} \\
Insecure randomness
    & No randomness consumed
    & \textcolor{passgreen}{\textbf{N/A}} \\
Storage collision / uninitialised storage
    & Independent contracts; \texttt{\^{}0.8.34} compiler check
    & \textcolor{passgreen}{\textbf{N/A}} \\
Short-address / calldata mis-encoding
    & EVM tooling mitigation in compiler
    & \textcolor{passgreen}{\textbf{N/A}} \\
Hidden fee / honeypot pattern
    & Symmetric ERC20; no asymmetric branch in code
    & \textcolor{passgreen}{\textbf{None present}} \\
Constructor / deployment front-run
    & Non-upgradeable; wiring atomic at construction
    & \textcolor{passgreen}{\textbf{No window}} \\
\bottomrule
\end{tabular}
\end{center}

In aggregate, the audit confirms that outside of the two findings
documented in Section~5 (C-1 and C-2), the protocol's defensive surface
behaves as designed against every attack vector exercised. The only
remaining recommendation is the gas-side optimisation in G-1
(Section~6.1).

\newpage

% ====================================================================
%  8. REMEDIATION CHECKLIST
% ====================================================================
\section{Remediation Checklist}

\begin{center}
\renewcommand{\arraystretch}{1.45}
\begin{tabular}{>{\bfseries}c p{8.8cm} >{\centering\arraybackslash}p{3cm}}
\toprule
\textbf{ID} & \textbf{Action} & \textbf{Type} \\
\midrule
C-1 & Add \texttt{minTokensOut} argument to \texttt{withdraw()} and
       switch \texttt{getSpotPrice()} to a TWAP over $\geq 30$ minutes;
       bound the mint with a DAO-tunable \texttt{MAX\_TOKENS\_PER\_USD}.
    & Code change \\
\midrule
C-2 & Introduce \texttt{MAX\_SUPPLY} on \texttt{i6token.sol} and check
       \texttt{totalSupply() + amount <= MAX\_SUPPLY} inside
       \texttt{mint()}.
    & Code change \\
\midrule
G-1 & DAO call: \texttt{system.setMaxDownlineDepth(600)}. Optionally
       lower the structural ceiling in the contract source on the next
       deployment.
    & Configuration \\
\bottomrule
\end{tabular}
\end{center}

\vspace{0.4cm}
After applying C-1 and C-2 we recommend re-running the
\texttt{25\_SpotPriceMint}, \texttt{26\_NoMaxSupply}, and
\texttt{40\_MEVSandwich} suites; each should flip from
\emph{demonstrating the vulnerability} to \emph{demonstrating the fix}.

\newpage

% ====================================================================
%  9. APPENDICES
% ====================================================================
\section{Appendix A -- Reproducing the Test Suite}

All tests live under
\texttt{test/Phase 4/Step 1/tests/} and inherit \texttt{BaseFork.t.sol}, which
pins the live BSC mainnet state.

\begin{lstlisting}[language=bash,caption={Running the suite},basicstyle=\ttfamily\footnotesize]
# Set any BSC RPC; an archive node is recommended for the gas suite.
export BSC_RPC_URL="https://bsc-rpc.publicnode.com"

# Optional reproducibility -- pin to a specific BSC block.
export BSC_BLOCK_NUMBER=0   # 0 means "latest"

# Run the full audit suite
forge test --match-path 'test/Phase 4/Step 1/tests/**/*.sol' -vv

# Run only the suites referenced in this report
forge test --match-path 'test/Phase 4/Step 1/tests/25_SpotPriceMint/*.sol' -vv
forge test --match-path 'test/Phase 4/Step 1/tests/26_NoMaxSupply/*.sol'   -vv
forge test --match-path 'test/Phase 4/Step 1/tests/38_GasAnalysis/*.sol'   -vv
forge test --match-path 'test/Phase 4/Step 1/tests/39_FlashLoan/*.sol'     -vv
forge test --match-path 'test/Phase 4/Step 1/tests/40_MEVSandwich/*.sol'   -vv
\end{lstlisting}

\section{Appendix B -- Disclaimer}

This audit report reflects the state of the codebase at the time of review and
the explicit attack surface enumerated in Section~\ref{sec:methodology}. No
audit can guarantee the absence of vulnerabilities. Findings outside the
listed scope, future code changes, off-chain key management, third-party
dependency behaviour, and economic / governance decisions taken by the DAO
multisig are not covered by this report. The recommendations in
Section~5 and Section~6 should be implemented, re-tested, and re-reviewed
before any reliance is placed on them.

\end{document}
